\documentclass[11pt,a4paper]{article}

\usepackage[margin=1in]{geometry}
\usepackage{times}
\usepackage{microtype}
\usepackage{amsmath,amssymb,amsfonts,amsthm}
\usepackage{booktabs}
\usepackage{multirow}
\usepackage{array}
\usepackage{adjustbox}
\usepackage{enumitem}
\usepackage{algorithm}
\usepackage{algorithmic}
\usepackage{graphicx}
\usepackage{subcaption}
\usepackage{float}
\usepackage{placeins}
\usepackage{xcolor}
\usepackage{url}
\usepackage[colorlinks=true,linkcolor=blue,citecolor=blue,urlcolor=blue]{hyperref}

\newtheorem{definition}{Definition}
\newtheorem{theorem}{Theorem}
\newtheorem{lemma}{Lemma}
\newtheorem{property}{Property}
\newtheorem{corollary}{Corollary}
\newtheorem{remark}{Remark}

\newcommand{\D}{\mathcal{D}}
\newcommand{\Vocab}{\mathcal{V}}
\newcommand{\Tokens}{\mathcal{X}}
\newcommand{\Txn}{\mathcal{T}}
\newcommand{\ByteStream}{\mathcal{S}}
\newcommand{\FP}{\operatorname{fp}}
\newcommand{\DIB}{\operatorname{dib}}
\newcommand{\home}{\operatorname{home}}
\newcommand{\toklen}{\operatorname{len}}
\newcommand{\tokid}{\operatorname{id}}
\newcommand{\rankop}{\operatorname{rank}}
\newcommand{\uniq}{\operatorname{uniq}}
\newcommand{\sort}{\operatorname{sort}}
\newcommand{\popcnt}{\operatorname{popcnt}}
\newcommand{\RSS}{\operatorname{RSS}}
\newcommand{\Throughput}{\operatorname{Throughput}}
\newcommand{\TokenRate}{\operatorname{TokenRate}}
\newcommand{\AvgProbe}{\operatorname{AvgProbe}}
\newcommand{\CollisionRate}{\operatorname{CollisionRate}}
\newcommand{\ExpansionOverhead}{\operatorname{ExpansionOverhead}}

\title{\textbf{FARO-Tokenizer: A Flat-Array Robin-Hood Offset Tokenizer for Zero-Dependency High-Performance Text Mining Pipelines}}

\author{
Quan Van\\
Independent Researcher\\
\texttt{vanhaminhquan2406@gmail.com}
}

\date{}

\begin{document}

\maketitle

\begin{abstract}
Massive text mining pipelines frequently spend substantial time transforming raw text into integer identifiers before downstream algorithms such as Apriori, FP-Growth, Eclat, or high-utility itemset mining can operate. Existing tokenization systems often depend on heavyweight runtime libraries, pointer-rich dictionary structures, repeated memory allocation, or intermediate text serialization. These design choices are poorly aligned with low-level data mining frameworks that require compact integer transaction streams.

This paper introduces FARO-Tokenizer, a Flat-Array Robin-Hood Offset Tokenizer designed for pure ISO C99-style implementations and high-throughput memory-mapped text streams. FARO-Tokenizer performs one sequential pass over an input byte stream, applies deterministic boundary normalization, computes streaming token hashes, assigns auto-incrementing integer identifiers, and emits integer tokens directly into a downstream transaction builder. Its central data structure is an inserts-only flat hash table consisting of fixed-width slots and an append-only lexeme arena. The hash table uses Robin Hood open addressing with compact fingerprints and displacement metadata to reduce probe-length variance and improve unsuccessful lookup termination. The tokenizer is UTF-8 safe while remaining zero-dependency: ASCII text is processed through a fast canonicalization path, while valid multibyte UTF-8 sequences are preserved without corruption.

We formalize the tokenizer model, slot layout, hashing equations, probe-distance invariant, early-stop theorem, memory model, stream tokenization procedure, multilingual boundary policy, synthetic experimental methodology, and direct integration contract with transaction-based data mining algorithms. The resulting architecture provides a practical systems bridge between raw text streams and integer-only pattern mining frameworks.
\end{abstract}

\noindent\textbf{Keywords:}
Tokenization; text mining; hash table; Robin Hood hashing; memory-mapped files; C99; UTF-8; itemset mining; FP-Growth; Apriori; SPMF.

\section{Introduction}

Text mining systems often begin with a deceptively expensive preprocessing stage: converting raw text into normalized tokens and mapping those tokens into compact integer identifiers. In high-performance data mining, this mapping is not merely a convenience. Algorithms such as Apriori, FP-Growth, Eclat, closed itemset mining, high-utility itemset mining, and stream mining operate most efficiently on integer item identifiers rather than variable-length strings \cite{AgrawalSrikant1994,HanPeiYin2000,Zaki2000,FournierViger2014SPMF}. If tokenization becomes memory-heavy or allocation-heavy, the downstream miner inherits unnecessary bottlenecks before the actual mining phase begins.

Traditional token dictionaries are often implemented using pointer-rich trees, chained hash tables, or general-purpose map containers. These structures are flexible, but they are not ideal for the workload considered in this paper. A tokenizer for offline mining has a simpler access pattern: it repeatedly performs get-or-create operations, assigns a new integer identifier only when a canonical token is first observed, and rarely needs deletion. This inserts-only dictionary workload allows more aggressive specialization than general-purpose string maps.

This paper proposes FARO-Tokenizer, a Flat-Array Robin-Hood Offset Tokenizer for massive memory-mapped text streams. FARO-Tokenizer is designed around four principles.

\begin{enumerate}[label=(\arabic*)]
\item \textbf{Flat memory layout}: the dictionary is a contiguous slot array, not a pointer graph.
\item \textbf{One-pass processing}: normalization, hashing, dictionary lookup, and ID emission are fused into one scan.
\item \textbf{Intern-once semantics}: each unique canonical token is stored exactly once in an append-only arena.
\item \textbf{Mining-native output}: emitted integer IDs are directly consumed by a transaction builder without intermediate string files.
\end{enumerate}

The proposed tokenizer is deliberately compatible with a low-level C99 implementation style. It does not require C++ containers, regular-expression libraries, ICU, Python runtime support, or external tokenization frameworks. The goal is not to compete with linguistically complete tokenizers; rather, it is to provide a deterministic, cache-friendly, zero-dependency tokenization layer for systems that prioritize throughput, reproducibility, and direct integration with data mining algorithms.

The contributions of this paper are as follows.

\begin{itemize}
\item We define the FARO tokenization model for memory-mapped text streams and integer-token mining pipelines.
\item We formalize a flat slot-array dictionary with compact metadata, displacement tracking, and append-only lexeme storage.
\item We define a streaming hash model based on byte-wise incremental hashing and length-aware finalization.
\item We prove the Robin Hood early-stop property used for efficient unsuccessful lookup.
\item We describe a zero-dependency multilingual boundary policy that is UTF-8 safe without requiring full Unicode segmentation.
\item We present tokenizer-to-miner contracts for line, document, sentence, and sliding-window transaction construction.
\item We propose a rigorous experimental methodology covering throughput, memory, collision behavior, probe distribution, rehash overhead, and downstream transaction compatibility.
\end{itemize}

\section{Background and Motivation}

\subsection{Text Tokenization for Mining}

Let a raw text corpus be represented as a byte stream
\[
S = \langle b_1,b_2,\dots,b_B\rangle,
\]
where each \(b_j\in\{0,\dots,255\}\). A tokenizer transforms this stream into a sequence of canonical tokens:
\[
X = \langle x_1,x_2,\dots,x_N\rangle,
\]
where each token \(x_i\) is a finite byte string after normalization and boundary handling.

For downstream data mining, the tokenizer must additionally construct a vocabulary map:
\[
\phi:\Tokens \rightarrow \mathbb{N}^+,
\]
where \(\Tokens\) is the set of observed canonical tokens and \(\phi(x)\) is the unique integer ID assigned to token \(x\). The emitted token stream is:
\[
Y=\langle \phi(x_1),\phi(x_2),\dots,\phi(x_N)\rangle.
\]

This integer stream can be grouped into transactions for itemset mining:
\[
\D=\{T_1,T_2,\dots,T_m\},
\qquad
T_j\subseteq \mathbb{N}^+.
\]

The critical systems problem is to compute \(\phi\) and \(Y\) with minimal memory traffic, minimal allocation, and minimal intermediate serialization.

\subsection{Why Flat Hashing Instead of Trees}

Trie-based tokenizers and pointer-based dictionaries can be useful when prefix lookup or language-specific segmentation is required. However, the mining-oriented tokenization task mainly requires exact canonical-token lookup. In such workloads, pointer-heavy dictionaries suffer from cache misses and allocator overhead. A flat open-addressed table keeps probe sequences contiguous and improves spatial locality \cite{AskitisZobel2005,Bender2021LinearProbing}.

Robin Hood hashing is especially suitable because it reduces probe-length variance by favoring keys that have traveled farther from their home bucket \cite{Celis1986RobinHood,Mitzenmacher2014RobinHood}. This matters in tokenization because high-frequency tokens repeatedly trigger successful lookups, while newly observed tokens require insertion. A stable and low-variance probing scheme improves predictability.

\subsection{Memory-Mapped Stream Processing}

For large files, memory mapping avoids explicit buffered read loops and allows the operating system to manage paging. A memory-mapped file can be scanned sequentially as an address range:
\[
S = [p,p+B).
\]
On POSIX-like systems, sequential-access hints may be used to inform the kernel of scan behavior \cite{LinuxMmap,LinuxMadvise,LinuxFadvise}. FARO-Tokenizer treats memory mapping as an ingestion strategy, while the core tokenization logic remains a single pass over bytes.

\section{The FARO Tokenizer Model}

\subsection{Canonical Token Function}

Let
\[
c:S^* \rightarrow S^*
\]
be a canonicalization function mapping a raw token byte span to a canonical token byte string. For ASCII text, \(c\) may apply case folding and punctuation trimming. For valid non-ASCII UTF-8 spans, FARO-Tokenizer preserves byte sequences unless a configured delimiter rule applies.

\begin{definition}[Canonical Token]
A canonical token is a non-empty byte string \(x\in\Tokens\) produced by applying boundary detection and canonicalization to a maximal token span in the input stream.
\end{definition}

\begin{definition}[Tokenizer Output]
Given byte stream \(S\), the tokenizer output is the integer sequence:
\[
Y(S)=\langle \phi(c(r_1)),\phi(c(r_2)),\dots,\phi(c(r_N))\rangle,
\]
where \(r_i\) are raw token spans and \(\phi\) is the auto-incrementing vocabulary map.
\end{definition}

\subsection{Design Requirements}

The design targets the following requirements.

\begin{enumerate}[label=(R\arabic*)]
\item \textbf{Zero dependency}: no runtime dependency on external tokenizer, Unicode, hash-table, or regular-expression libraries.
\item \textbf{One pass}: every input byte is inspected \(O(1)\) times.
\item \textbf{No per-token allocation}: duplicate tokens do not allocate memory.
\item \textbf{Cache-friendly dictionary}: dictionary probes follow contiguous memory.
\item \textbf{Stable ID assignment}: the first occurrence of a token determines its integer ID.
\item \textbf{UTF-8 safety}: valid multibyte sequences are not split or corrupted.
\item \textbf{Mining-native emission}: output IDs can be consumed by itemset miners without string reconstruction.
\end{enumerate}

\section{Flat-Array Dictionary Architecture}

\subsection{Slot Array}

The dictionary is a table:
\[
A=\langle a_0,a_1,\dots,a_{C-1}\rangle,
\]
where \(C=2^k\) is the table capacity. Each slot stores:
\[
a_s = \langle m_s,o_s,\ell_s,z_s\rangle,
\]
where \(m_s\) is metadata, \(o_s\) is an offset into the lexeme arena, \(\ell_s\) is token length, and \(z_s\) is the assigned integer ID.

\begin{table}[H]
\centering
\caption{FARO dictionary slot layout.}
\label{tab:slot-layout}
\small
\begin{adjustbox}{max width=\textwidth}
\begin{tabular}{llll}
\toprule
Field & Symbol & Logical Width & Meaning \\
\midrule
Metadata & \(m_s\) & 32 bits & fingerprint, displacement, and flags \\
Arena offset & \(o_s\) & 32 bits & byte offset of canonical token in arena \\
Token length & \(\ell_s\) & 32 bits & token length in bytes \\
Token ID & \(z_s\) & 32 bits & auto-incremented integer identifier \\
\bottomrule
\end{tabular}
\end{adjustbox}
\end{table}

The metadata word is interpreted as:
\[
m_s = \FP_s + 2^8\DIB_s + 2^{24}F_s,
\]
where \(\FP_s\in[0,255]\) is an 8-bit fingerprint, \(\DIB_s\in[0,2^{16}-1]\) is displacement from the initial bucket, and \(F_s\) stores optional flags. The fingerprint value \(0\) is reserved for empty slots.

\subsection{Lexeme Arena}

All unique canonical tokens are stored in an append-only byte arena:
\[
R = \langle r_0,r_1,\dots,r_{A-1}\rangle.
\]
A slot \(\langle m_s,o_s,\ell_s,z_s\rangle\) refers to token bytes:
\[
R[o_s:o_s+\ell_s).
\]

\begin{definition}[Interned Token]
A token \(x\) is interned if exactly one byte copy of \(x\) is stored in the lexeme arena and all future occurrences reuse its slot ID.
\end{definition}

This model separates hot lookup metadata from cold variable-length token bytes. During probing, most slots are rejected using fingerprint and length before byte comparison.

\subsection{Borrowed-View Mode}

FARO-Tokenizer also admits a borrowed-view mode where \(o_s\) is interpreted as an offset into the input mapping rather than the lexeme arena. This mode is correct only when:
\[
c(r)=r,
\]
and the mapped byte stream remains valid for the lifetime of the dictionary. Borrowed-view mode is therefore a specialized optimization for identity canonicalization and single-mapping pipelines.

\section{Hashing Model}

\subsection{Streaming Byte Hash}

Let a canonical token be:
\[
x=\langle b_1,b_2,\dots,b_\ell\rangle.
\]
FARO-Tokenizer computes a streaming hash while building the token:
\[
H_0 = \eta,
\]
\[
H_j = (H_{j-1}\oplus b_j)\cdot p \pmod{2^{64}},
\qquad 1\le j\le \ell,
\]
where \(\eta\) is the 64-bit offset basis and \(p\) is the 64-bit FNV prime \cite{Eastlake2026FNV}.

At token end, the length is mixed:
\[
G(x)=H_\ell\oplus \ell,
\]
and an avalanche finalizer is applied:
\[
h(x)=\operatorname{fmix64}(G(x)).
\]
The home bucket is:
\[
\home(x)=h(x)\mathbin{\&}(C-1).
\]

\subsection{Fingerprint}

The compact slot fingerprint is:
\[
\FP(x)=
\begin{cases}
1, & \text{if } (h(x)\gg 56)=0,\\
h(x)\gg 56, & \text{otherwise}.
\end{cases}
\]
This reserves fingerprint zero for empty slots.

\begin{property}[Fingerprint Rejection]
If two tokens \(x\) and \(y\) satisfy \(\FP(x)\ne\FP(y)\), then they cannot match in the dictionary and byte comparison is unnecessary.
\end{property}

\section{Robin Hood Open Addressing}

\subsection{Probe Distance}

Let \(x\) be stored in slot \(s\). Its distance to initial bucket is:
\[
\delta(x,s)=(s-\home(x))\bmod C.
\]

The table load factor is:
\[
\alpha=\frac{n}{C},
\]
where \(n\) is the number of occupied slots.

\subsection{Robin Hood Invariant}

Robin Hood insertion maintains the principle that an incoming key with a larger displacement may swap with an occupant having a smaller displacement. Informally, keys that are farther from home are allowed to steal positions from keys closer to home.

\begin{definition}[Robin Hood Validity]
A hash table state is Robin-Hood-valid if every insertion has followed the rule:
\[
d_{\text{incoming}}>d_{\text{occupant}}
\quad\Rightarrow\quad
\text{swap incoming and occupant}.
\]
\end{definition}

\begin{theorem}[Early-Stop Property]
Assume a Robin-Hood-valid table. During lookup of token \(x\), suppose the probe reaches slot \(s\) with current displacement \(d\). If slot \(s\) is occupied by token \(y\) with:
\[
\delta(y,s)<d,
\]
then \(x\) is not present in the table.
\end{theorem}

\begin{proof}
Assume for contradiction that \(x\) is present beyond slot \(s\). When \(x\) was inserted, it must have reached slot \(s\) with displacement \(d\). At that moment, the occupant \(y\) had displacement strictly smaller than \(d\). The Robin Hood rule would have swapped \(x\) ahead of \(y\). Therefore \(x\) could not have remained behind \(y\). This contradicts the assumption that \(x\) is present beyond \(s\). Hence \(x\) is absent.
\end{proof}

\begin{corollary}[Unsuccessful Lookup Termination]
Unsuccessful lookup in a Robin-Hood-valid FARO dictionary can terminate before reaching an empty slot whenever the current displacement exceeds the stored displacement of the examined slot.
\end{corollary}

\section{Dictionary Operations}

\subsection{Get-or-Create Semantics}

The central operation is:
\[
\operatorname{getOrCreate}(x)=
\begin{cases}
\phi(x), & \text{if }x\in\Tokens,\\
|\Tokens|+1, & \text{otherwise, after inserting }x.
\end{cases}
\]

The operation is inserts-only. Deletion is not required because the tokenizer builds a vocabulary monotonically during a corpus scan.

\subsection{Mathematical Algorithm}

\begin{algorithm}[H]
\caption{FARO Get-or-Create}
\label{alg:get-or-create}
\begin{algorithmic}[1]
\REQUIRE canonical token \(x\), hash \(h(x)\), table \(A\), arena \(R\)
\ENSURE integer token identifier \(\phi(x)\)
\STATE \(s\leftarrow \home(x)\)
\STATE \(d\leftarrow 0\)
\STATE \(f\leftarrow \FP(x)\)
\STATE candidate token state remains unmaterialized
\LOOP
    \IF{slot \(A_s\) is empty}
        \STATE materialize \(x\) into arena \(R\) if not already materialized
        \STATE assign next integer identifier to \(x\)
        \STATE store \(\langle f,d,o_x,\toklen(x),\phi(x)\rangle\) in \(A_s\)
        \STATE \RETURN \(\phi(x)\)
    \ENDIF
    \IF{\(\FP(A_s)=f\), \(\ell_s=\toklen(x)\), and arena bytes equal \(x\)}
        \STATE \RETURN \(z_s\)
    \ENDIF
    \IF{\(\DIB(A_s)<d\)}
        \STATE materialize \(x\) into arena \(R\) if not already materialized
        \STATE swap candidate state with occupant of \(A_s\)
        \STATE update candidate fingerprint and displacement
    \ENDIF
    \STATE \(s\leftarrow (s+1)\bmod C\)
    \STATE \(d\leftarrow d+1\)
\ENDLOOP
\end{algorithmic}
\end{algorithm}

\subsection{Late Materialization}

\begin{definition}[Late Materialization]
A token occurrence is late-materialized if its canonical bytes are copied into the arena only after the lookup process determines that insertion is unavoidable.
\end{definition}

\begin{property}[Duplicate-Copy Avoidance]
If token \(x\) already exists in the dictionary, late materialization avoids copying \(x\) into the arena.
\end{property}

This is important for natural language corpora, where token frequencies are highly skewed and common words appear many times.

\section{Stream Tokenization}

\subsection{Tokenizer State Machine}

FARO-Tokenizer scans the byte stream using a deterministic state machine:
\[
Q=\{\textsf{Outside},\textsf{InsideASCII},\textsf{InsideUTF8},\textsf{ErrorRecovery}\}.
\]

At each byte position \(j\), the tokenizer classifies \(b_j\) as one of:
\[
\textsf{Word},\quad
\textsf{Delimiter},\quad
\textsf{Joiner},\quad
\textsf{UTF8Lead},\quad
\textsf{UTF8Continuation},\quad
\textsf{Invalid}.
\]

The scanner maintains a scratch token buffer \(W\), current length \(\ell\), and streaming hash state \(H_\ell\). When a token boundary is reached, the tokenizer finalizes the hash, invokes get-or-create, and emits the integer ID.

\begin{algorithm}[H]
\caption{FARO Tokenize Stream}
\label{alg:tokenize-stream}
\begin{algorithmic}[1]
\REQUIRE byte stream \(S=\langle b_1,\dots,b_B\rangle\), dictionary \(D\), emitter \(E\)
\ENSURE integer token stream
\STATE initialize token buffer \(W\leftarrow\emptyset\)
\STATE initialize state \(q\leftarrow\textsf{Outside}\)
\STATE initialize streaming hash \(H\leftarrow\eta\)
\FOR{\(j=1\) to \(B\)}
    \STATE classify byte \(b_j\)
    \IF{\(b_j\) is ASCII word byte}
        \STATE fold \(b_j\) if ASCII case folding is enabled
        \STATE append canonical byte to \(W\)
        \STATE update streaming hash
        \STATE \(q\leftarrow\textsf{InsideASCII}\)
    \ELSIF{\(b_j\) is valid UTF-8 lead byte}
        \STATE validate full UTF-8 sequence beginning at \(j\)
        \IF{sequence is valid and not configured as delimiter}
            \STATE append all bytes of the sequence to \(W\)
            \STATE update streaming hash with each byte
            \STATE advance \(j\) to the end of the sequence
            \STATE \(q\leftarrow\textsf{InsideUTF8}\)
        \ELSE
            \STATE finalize current token if \(W\ne\emptyset\)
            \STATE \(q\leftarrow\textsf{Outside}\)
        \ENDIF
    \ELSIF{\(b_j\) is an internal joiner and neighboring bytes are token bytes}
        \STATE append joiner byte to \(W\)
        \STATE update streaming hash
    \ELSE
        \IF{\(W\ne\emptyset\)}
            \STATE finalize hash of \(W\)
            \STATE \(z\leftarrow \operatorname{getOrCreate}(W)\)
            \STATE emit \(z\) to downstream consumer
            \STATE reset \(W\leftarrow\emptyset\)
        \ENDIF
        \STATE \(q\leftarrow\textsf{Outside}\)
    \ENDIF
\ENDFOR
\IF{\(W\ne\emptyset\)}
    \STATE finalize hash of \(W\)
    \STATE emit \(\operatorname{getOrCreate}(W)\)
\ENDIF
\end{algorithmic}
\end{algorithm}

\subsection{Complexity of Tokenization}

Let \(B\) be the number of input bytes and \(N\) the number of emitted tokens. Let \(P_i\) be the number of probes for token \(i\). The total time is:
\[
T(B,N)=\Theta(B)+\Theta\left(\sum_{i=1}^{N}P_i\right).
\]
Under a controlled load factor \(\alpha<\alpha_{\max}\), expected probing is constant:
\[
\mathbb{E}[P_i]=O(1),
\]
so:
\[
\mathbb{E}[T(B,N)]=\Theta(B+N).
\]
Since \(N\le B\), the expected end-to-end runtime is:
\[
\mathbb{E}[T(B,N)]=\Theta(B).
\]

\section{Multilingual and Boundary Handling}

\subsection{ASCII Fast Path}

For bytes \(b<128\), FARO-Tokenizer uses a constant-time table:
\[
C_{\text{ascii}}:\{0,\dots,127\}\rightarrow\{\textsf{Word},\textsf{Delimiter},\textsf{Joiner},\textsf{Control}\}.
\]
A second table implements ASCII folding:
\[
F_{\text{ascii}}(b)=
\begin{cases}
b+32, & \text{if }65\le b\le 90,\\
b, & \text{otherwise}.
\end{cases}
\]

\subsection{Joiner Rule}

Let \(J=\{\texttt{-},\texttt{\_},\texttt{'}\}\) be the ASCII joiner set. A joiner byte at position \(j\) is retained inside a token only if:
\[
\operatorname{isTokenByte}(b_{j-1})=1
\quad\text{and}\quad
\operatorname{isTokenByte}(b_{j+1})=1.
\]
Otherwise it is treated as a delimiter.

This rule keeps forms such as hyphenated terms and contractions while stripping leading and trailing punctuation.

\subsection{UTF-8 Safety}

FARO-Tokenizer follows a UTF-8-safe rather than fully Unicode-complete design. UTF-8 lead bytes determine sequence length, and continuation bytes must satisfy:
\[
b \in [128,191].
\]
A byte sequence is valid only if the lead byte and its required continuation bytes form a legal UTF-8 unit \cite{Yergeau2003UTF8}.

\begin{definition}[UTF-8 Safe Tokenization]
A tokenizer is UTF-8 safe if it never splits a valid multibyte UTF-8 sequence into separate corrupted byte fragments and never treats continuation bytes as independent ASCII tokens.
\end{definition}

\begin{property}[ASCII Compatibility]
Since UTF-8 preserves the ASCII range, the ASCII fast path and UTF-8 validation path can coexist without ambiguity.
\end{property}

\subsection{Unicode Scope}

Full Unicode word segmentation and full Unicode case folding are not assumed. Unicode text segmentation is rule-based and language-sensitive \cite{UnicodeUAX29}. Full Unicode case folding may expand strings and may not preserve normalization \cite{UnicodeCaseFolding}. Therefore, FARO-Tokenizer adopts the following zero-dependency policy.

\begin{enumerate}[label=(\arabic*)]
\item ASCII case folding is supported by default.
\item Valid non-ASCII UTF-8 sequences are preserved as token bytes.
\item A small configurable delimiter set may recognize common non-ASCII punctuation.
\item Full Unicode segmentation is treated as an optional external preprocessing layer, not a core dependency.
\end{enumerate}

\section{Memory Model}

Let \(U\) be the number of unique canonical tokens, \(C\) be the hash table capacity, and \(A\) be the number of bytes used by the lexeme arena. If each slot occupies \(S_{\text{slot}}\) bytes, heap memory is approximately:
\[
M_{\text{heap}}=S_{\text{slot}}C+A+M_{\text{scratch}}+M_{\text{emit}}.
\]
With a 16-byte slot:
\[
M_{\text{heap}}\approx 16C+A+M_{\text{scratch}}+M_{\text{emit}}.
\]

If the maximum load factor is \(\alpha_{\max}\), then:
\[
C \ge \left\lceil \frac{U}{\alpha_{\max}}\right\rceil_2,
\]
where \(\lceil\cdot\rceil_2\) denotes rounding up to the next power of two.

Thus:
\[
M_{\text{heap}}\approx
16\left\lceil \frac{U}{\alpha_{\max}}\right\rceil_2
+
\sum_{x\in\Tokens}\toklen(x)
+
O(1)
+
M_{\text{emit}}.
\]

If input is memory-mapped, process RSS may include file-backed pages:
\[
M_{\RSS}=M_{\text{heap}}+M_{\text{file-backed}}+M_{\text{runtime}},
\]
so experiments must distinguish heap-owned memory from resident mapped pages.

\section{Transaction Construction for Data Mining}

\subsection{Token Stream to Transactions}

The emitted ID stream:
\[
Y=\langle y_1,y_2,\dots,y_N\rangle
\]
must be grouped into transactions:
\[
\D_Y=\{T_1,T_2,\dots,T_m\}.
\]

FARO-Tokenizer supports three transaction modes.

\begin{definition}[Document Transaction Mode]
Each logical document, line, record, or paragraph becomes one transaction:
\[
T_j=\uniq(\{y_i : y_i \text{ occurs in document }j\}).
\]
\end{definition}

\begin{definition}[Sentence Transaction Mode]
Each detected sentence span becomes one transaction:
\[
T_j=\uniq(\{y_i : y_i \text{ occurs in sentence }j\}).
\]
\end{definition}

\begin{definition}[Sliding-Window Transaction Mode]
Given window size \(w\) and stride \(s\), define:
\[
T_j=
\uniq\left(
\{y_{1+(j-1)s},\dots,y_{\min(N,(j-1)s+w)}\}
\right).
\]
\end{definition}

\subsection{Set Semantics}

Classical itemset mining transactions are sets. Therefore, every transaction emitted to a miner must satisfy:
\[
T_j=\sort(\uniq(T_j)).
\]
This removes duplicate token IDs within the same transaction and imposes a deterministic total order.

\subsection{SPMF-Compatible Contract}

A transaction is SPMF-compatible if:
\[
T_j=\langle t_{j1},t_{j2},\dots,t_{jk}\rangle,
\]
where:
\[
1\le t_{j1}<t_{j2}<\cdots<t_{jk}.
\]

Thus, FARO-Tokenizer can feed miners in two ways.

\begin{enumerate}[label=(\arabic*)]
\item \textbf{In-memory contract}: each transaction is passed directly to a mining callback as an array of sorted unique integer IDs.
\item \textbf{Text compatibility contract}: each transaction is serialized as one line of space-separated positive integers.
\end{enumerate}

The in-memory contract avoids intermediate disk I/O and is the preferred path for embedded C mining frameworks.

\section{Experimental Methodology}

\subsection{Research Questions}

The evaluation should answer the following questions.

\begin{itemize}
\item RQ1: Does FARO-Tokenizer achieve higher throughput than pointer-heavy dictionary designs?
\item RQ2: Does Robin Hood probing reduce average, maximum, and tail probe lengths compared with linear probing?
\item RQ3: What is the memory cost per unique token under different load factors?
\item RQ4: How much overhead is caused by table expansion and rehashing?
\item RQ5: Does UTF-8-safe scanning preserve correctness while maintaining high throughput?
\item RQ6: Can the tokenizer feed itemset miners without intermediate string materialization?
\end{itemize}

\subsection{Benchmark Variants}

The following variants should be evaluated.

\begin{table}[H]
\centering
\caption{Tokenizer variants for ablation study.}
\label{tab:variants}
\small
\begin{adjustbox}{max width=\textwidth}
\begin{tabular}{lll}
\toprule
Variant & Collision Strategy & Purpose \\
\midrule
LP-Raw & linear probing without fingerprint & baseline open addressing \\
LP-FP & linear probing with fingerprint & isolates fingerprint rejection \\
RH-FP & Robin Hood with fingerprint & isolates displacement balancing \\
RH-Arena & Robin Hood with arena interning & default normalized tokenizer \\
RH-Borrow & Robin Hood with borrowed views & zero-copy identity-canonicalization mode \\
\bottomrule
\end{tabular}
\end{adjustbox}
\end{table}

\subsection{Datasets}

A rigorous evaluation should include both real and synthetic datasets.

\begin{enumerate}[label=(\arabic*)]
\item English technical documents.
\item Multilingual web text.
\item Code-switched social-media-like text.
\item Log files with repeated templates and rare identifiers.
\item Synthetic Zipfian corpora of 500 MB and 1 GB.
\item Downstream transaction streams generated from sliding windows.
\end{enumerate}

Synthetic text should approximate natural language through a Zipf-Mandelbrot token frequency distribution:
\[
P(r)=\frac{1}{Z}\cdot\frac{1}{(r+q)^s},
\]
where \(r\) is token rank, \(s>1\) controls skew, \(q\ge0\) is a shift parameter, and:
\[
Z=\sum_{r=1}^{V}\frac{1}{(r+q)^s}.
\]
This reflects the heavy-tailed vocabulary behavior observed in natural language \cite{Piantadosi2014Zipf}.

\subsection{Metrics}

Throughput is:
\[
\Throughput_{MiB/s}=
\frac{B}{2^{20}\cdot t},
\]
where \(B\) is bytes processed and \(t\) is elapsed seconds.

Token rate is:
\[
\TokenRate=
\frac{N}{t},
\]
where \(N\) is the number of emitted tokens.

Average probe length is:
\[
\AvgProbe=
\frac{\sum_{i=1}^{Q}p_i}{Q},
\]
where \(Q\) is the number of dictionary queries and \(p_i\) is the probe count of query \(i\).

Collision rate is:
\[
\CollisionRate=
\frac{Q_{\text{collision}}}{Q}.
\]

Expansion overhead is:
\[
\ExpansionOverhead=
\frac{t_{\text{rehash}}}{t_{\text{total}}}.
\]

Memory per unique token is:
\[
M_{\text{per-token}}=
\frac{M_{\text{slots}}+M_{\text{arena}}}{U}.
\]

Probe tail behavior should report:
\[
P_{95},\quad P_{99},\quad P_{\max},
\]
where \(P_p\) is the \(p\)-th percentile probe length.

\subsection{Measurement Protocol}

Timing should use a monotonic clock:
\[
t=t_{\text{end}}-t_{\text{start}},
\]
where both timestamps come from a monotonic time source. Process-level memory should report peak resident usage and should distinguish heap allocation from memory-mapped file residency \cite{LinuxClockGettime,LinuxGetrusage,LinuxProcStatus}.

Each benchmark configuration should be executed at least five times:
\[
R=\{r_1,r_2,r_3,r_4,r_5\}.
\]
The reported center should be the median:
\[
\operatorname{median}(R),
\]
and variability should be the interquartile range:
\[
IQR(R)=Q_3(R)-Q_1(R).
\]

Cold-cache and warm-cache measurements should be reported separately because memory-mapped file scans can differ depending on page-cache state.

\subsection{Expected Experimental Tables}

\begin{table}[H]
\centering
\caption{Core performance metrics to report.}
\label{tab:metrics}
\small
\begin{adjustbox}{max width=\textwidth}
\begin{tabular}{llllllll}
\toprule
Dataset & Variant & MiB/s & Tokens/s & Unique & AvgProbe & MaxProbe & Peak RSS \\
\midrule
Corpus-A & LP-Raw & -- & -- & -- & -- & -- & -- \\
Corpus-A & LP-FP & -- & -- & -- & -- & -- & -- \\
Corpus-A & RH-FP & -- & -- & -- & -- & -- & -- \\
Corpus-A & RH-Arena & -- & -- & -- & -- & -- & -- \\
Corpus-A & RH-Borrow & -- & -- & -- & -- & -- & -- \\
\bottomrule
\end{tabular}
\end{adjustbox}
\end{table}

\begin{table}[H]
\centering
\caption{Expansion and memory behavior.}
\label{tab:expansion-memory}
\small
\begin{adjustbox}{max width=\textwidth}
\begin{tabular}{lllllll}
\toprule
Dataset & Variant & Load Factor & Rehash Count & Rehash Time Ratio & Slot MB & Arena MB \\
\midrule
Corpus-A & RH-Arena & -- & -- & -- & -- & -- \\
Corpus-B & RH-Arena & -- & -- & -- & -- & -- \\
Corpus-C & RH-Arena & -- & -- & -- & -- & -- \\
\bottomrule
\end{tabular}
\end{adjustbox}
\end{table}

\subsection{Recommended Visualizations}

The paper should include:

\begin{itemize}
\item throughput by tokenizer variant,
\item average probe length versus load factor,
\item \(P_{95}\) and \(P_{99}\) probe length versus load factor,
\item memory per unique token by variant,
\item rehash overhead by dataset size,
\item downstream transaction throughput,
\item cold-cache versus warm-cache throughput.
\end{itemize}

\section{Theoretical Evaluation}

\subsection{Time Complexity}

Let \(B\) be input bytes, \(N\) emitted tokens, \(U\) unique tokens, and \(P_i\) probe count for token \(i\). FARO-Tokenizer performs:
\[
T(B,N)=\Theta(B)+\Theta\left(\sum_{i=1}^{N}P_i\right).
\]
Under controlled load factor and non-adversarial hashing:
\[
\mathbb{E}[P_i]=O(1),
\]
so:
\[
\mathbb{E}[T(B,N)]=\Theta(B+N)=\Theta(B).
\]

\subsection{Memory Complexity}

With slot size \(S_{\text{slot}}\), capacity \(C\), arena byte count \(A\), and transaction buffer memory \(M_{\text{txn}}\):
\[
M=O(S_{\text{slot}}C+A+M_{\text{txn}}).
\]
Since:
\[
C=O\left(\frac{U}{\alpha_{\max}}\right),
\]
we obtain:
\[
M=O(U+A+M_{\text{txn}}).
\]
If average canonical token length is \(\bar{\ell}\), then:
\[
A\approx U\bar{\ell},
\]
and:
\[
M=O(U\bar{\ell}+U+M_{\text{txn}}).
\]

\subsection{Output Sensitivity}

The tokenizer is output-sensitive with respect to emitted token count and vocabulary size:
\[
T=\Theta(B)+O(N\cdot \bar{P}),
\]
\[
M=O(U\bar{\ell}+C).
\]
Thus, repeated tokens increase runtime through lookups but do not increase dictionary storage.

\section{Correctness}

\begin{theorem}[Vocabulary Uniqueness]
For every canonical token \(x\), FARO-Tokenizer assigns exactly one integer ID \(\phi(x)\).
\end{theorem}

\begin{proof}
If \(x\) is absent, get-or-create inserts \(x\) into exactly one empty or swapped slot and assigns the next unused integer ID. If \(x\) is present, the dictionary lookup finds a slot with matching fingerprint, length, and byte sequence, and returns the existing ID. Since insertion occurs only when no existing equivalent token is found, no two IDs are assigned to the same canonical token.
\end{proof}

\begin{theorem}[Token Stream Determinism]
Given the same input byte stream, canonicalization policy, hash function, initial capacity policy, and transaction boundary policy, FARO-Tokenizer emits the same integer sequence across runs.
\end{theorem}

\begin{proof}
The byte scan order is deterministic. Canonicalization is deterministic. The hash function and probing rule are deterministic. IDs are assigned in first-occurrence order. Therefore the emitted integer sequence is deterministic.
\end{proof}

\begin{theorem}[Transaction Set Correctness]
For any transaction span \(S_j\), the transaction emitted under set semantics is:
\[
T_j=\sort(\uniq(\{\phi(x):x\text{ occurs in }S_j\})).
\]
\end{theorem}

\begin{proof}
The transaction builder accumulates token IDs from the span, sorts them under the configured total order, and removes duplicates. Therefore each emitted transaction is exactly the sorted unique set of token IDs occurring in the span.
\end{proof}

\section{Discussion}

FARO-Tokenizer is not intended to replace full linguistic tokenizers. Its purpose is narrower and systems-oriented: it converts massive raw text streams into deterministic integer streams suitable for mining. This design choice explains why full Unicode segmentation is outside the core architecture. Full segmentation requires extensive rule tables and language-specific behavior, while FARO-Tokenizer prioritizes pure C99 portability, predictable memory, and direct integration with low-level miners.

The architecture is especially appropriate for data mining tasks where token identity and co-occurrence matter more than linguistically perfect token boundaries. Examples include log mining, keyword co-occurrence mining, document transaction construction, sliding-window itemset mining, and high-throughput preprocessing for integer-only mining algorithms.

\section{Limitations}

FARO-Tokenizer has several limitations.

First, it is not a full Unicode word segmenter. It is UTF-8 safe but does not implement the full Unicode text segmentation algorithm.

Second, ASCII folding is simple and deterministic but does not cover full Unicode case folding.

Third, Robin Hood hashing assumes a non-adversarial workload. If the input is adversarially constructed to collide under the hash function, performance may degrade.

Fourth, borrowed-view mode is safe only when the input mapping remains valid and canonicalization is identity.

Fifth, ID assignment is first-occurrence based. If distributed tokenization is required, an additional vocabulary merge protocol is needed.

\section{Conclusion}

This paper introduced FARO-Tokenizer, a flat-array Robin-Hood offset tokenizer for zero-dependency high-performance text mining pipelines. FARO-Tokenizer performs one-pass tokenization over memory-mapped text streams, uses streaming hash computation, stores unique canonical tokens exactly once, assigns auto-incrementing integer IDs, and emits mining-ready integer transactions.

The proposed architecture replaces pointer-heavy dictionaries with a flat contiguous hash table and append-only lexeme arena. Its Robin Hood probing model reduces probe variance and enables early unsuccessful lookup termination. Its UTF-8-safe boundary policy preserves multibyte text without requiring external Unicode libraries. Its output contract allows direct integration with Apriori, FP-Growth, Eclat, and SPMF-style transaction formats.

Future work includes parallel vocabulary construction, NUMA-aware table partitioning, adversarial hash hardening, optional generated Unicode case-fold tables, streaming compression of emitted ID sequences, and direct coupling with high-utility and high-occupancy itemset miners.

\bibliographystyle{plain}
\begin{thebibliography}{99}

\bibitem{AgrawalSrikant1994}
Rakesh Agrawal and Ramakrishnan Srikant.
\newblock Fast algorithms for mining association rules.
\newblock In \emph{Proceedings of the 20th International Conference on Very Large Data Bases}, pages 487--499, 1994.

\bibitem{HanPeiYin2000}
Jiawei Han, Jian Pei, and Yiwen Yin.
\newblock Mining frequent patterns without candidate generation.
\newblock In \emph{Proceedings of the ACM SIGMOD International Conference on Management of Data}, pages 1--12, 2000.

\bibitem{Zaki2000}
Mohammed J. Zaki.
\newblock Scalable algorithms for association mining.
\newblock \emph{IEEE Transactions on Knowledge and Data Engineering}, 12(3):372--390, 2000.

\bibitem{FournierViger2014SPMF}
Philippe Fournier-Viger, Antonio Gomariz, Ted Gueniche, Azadeh Soltani, Cheng-Wei Wu, and Vincent S. Tseng.
\newblock SPMF: A Java open-source pattern mining library.
\newblock \emph{Journal of Machine Learning Research}, 15:3389--3393, 2014.

\bibitem{AskitisZobel2005}
Nikolas Askitis and Justin Zobel.
\newblock Cache-conscious collision resolution in string hash tables.
\newblock In \emph{Proceedings of the International Symposium on String Processing and Information Retrieval}, pages 91--102, 2005.

\bibitem{Bender2021LinearProbing}
Michael A. Bender, Bradley C. Kuszmaul, and William Kuszmaul.
\newblock Linear probing revisited.
\newblock \emph{arXiv preprint arXiv:2107.01250}, 2021.

\bibitem{Celis1986RobinHood}
Pedro Celis.
\newblock Robin Hood hashing.
\newblock PhD thesis, University of Waterloo, 1986.

\bibitem{Mitzenmacher2014RobinHood}
Michael Mitzenmacher.
\newblock A new approach to analyzing Robin Hood hashing.
\newblock \emph{arXiv preprint arXiv:1401.7616}, 2014.

\bibitem{Eastlake2026FNV}
Donald E. Eastlake, Glenn Fowler, Landon Curt Noll, Kiem-Phong Vo, and Tony Hansen.
\newblock The FNV non-cryptographic hash algorithm.
\newblock RFC 9923, Internet Research Task Force, 2026.

\bibitem{Yergeau2003UTF8}
Fran\c{c}ois Yergeau.
\newblock UTF-8, a transformation format of ISO 10646.
\newblock RFC 3629, Internet Engineering Task Force, 2003.

\bibitem{UnicodeUAX29}
Unicode Consortium.
\newblock Unicode text segmentation.
\newblock Unicode Standard Annex \#29.

\bibitem{UnicodeCaseFolding}
Unicode Consortium.
\newblock Case folding.
\newblock Unicode Character Database.

\bibitem{LinuxMmap}
Michael Kerrisk.
\newblock \texttt{mmap(2)}: Linux manual page.
\newblock Linux man-pages project.

\bibitem{LinuxMadvise}
Michael Kerrisk.
\newblock \texttt{madvise(2)}: Linux manual page.
\newblock Linux man-pages project.

\bibitem{LinuxFadvise}
Michael Kerrisk.
\newblock \texttt{posix\_fadvise(2)}: Linux manual page.
\newblock Linux man-pages project.

\bibitem{LinuxClockGettime}
Michael Kerrisk.
\newblock \texttt{clock\_gettime(3)}: Linux manual page.
\newblock Linux man-pages project.

\bibitem{LinuxGetrusage}
Michael Kerrisk.
\newblock \texttt{getrusage(2)}: Linux manual page.
\newblock Linux man-pages project.

\bibitem{LinuxProcStatus}
Michael Kerrisk.
\newblock \texttt{proc\_pid\_status(5)}: Linux manual page.
\newblock Linux man-pages project.

\bibitem{Piantadosi2014Zipf}
Steven T. Piantadosi.
\newblock Zipf's word frequency law in natural language: A critical review and future directions.
\newblock \emph{Psychonomic Bulletin \& Review}, 21:1112--1130, 2014.

\bibitem{ApplebyMurmurHash}
Austin Appleby.
\newblock MurmurHash3.
\newblock SMHasher public domain reference implementation.

\bibitem{Pagh2006LinearProbing}
Anna Pagh, Rasmus Pagh, and Milan Ru\v{z}i\'c.
\newblock Linear probing with constant independence.
\newblock \emph{SIAM Journal on Computing}, 39(3):1107--1120, 2009.

\bibitem{DillingerManolios2004Bloom}
Peter C. Dillinger and Panagiotis Manolios.
\newblock Bloom filters in probabilistic verification.
\newblock In \emph{Proceedings of the International Conference on Formal Methods in Computer-Aided Design}, pages 367--381, 2004.

\bibitem{MoffatTurpin2002Compression}
Alistair Moffat and Andrew Turpin.
\newblock Compression and coding algorithms.
\newblock Kluwer Academic Publishers, 2002.

\bibitem{WittenMoffatBell1999Managing}
Ian H. Witten, Alistair Moffat, and Timothy C. Bell.
\newblock Managing Gigabytes: Compressing and Indexing Documents and Images.
\newblock Morgan Kaufmann, 1999.

\bibitem{ManningRaghavanSchutze2008IR}
Christopher D. Manning, Prabhakar Raghavan, and Hinrich Sch\"utze.
\newblock Introduction to Information Retrieval.
\newblock Cambridge University Press, 2008.

\end{thebibliography}

\end{document}
